\section{Introduction}
In this chapter, value function approximation methods will be covered. Unlike the tabular methods, the approximate method tries to
approximate the state value or action value of state-action pairs, using a function. This solves the curse of dimensionality, which is a critical problem for tabular methods.

But there are tradeoffs, too. For example, the convergence property breaks down easily. When off-policy, function approximation, and bootstrapping meets together,
(\textbf{the deadly triad}, \cite{Sutton2020}), the algorithm lost its convergence property and fails to find optimal policy. There are various ways to solve this problem.
Some solve this with practical methods (replay buffer, PER, double learning, etc.), and some solve this theoretically (GTD2, Emphatic-TD, etc.).
I'll try to cover both of them in this note.

\section{TD Learning of State Values based on Function Approximation}
Methods of estimating the state value with a function approximation is introduced in this section.
    \subsection{Objective Function}
        Let $v_\pi(s)$ be the true state value function and let $\hat{v}_\pi(s, w)$ be the approximation of it with a parameter $w$.
        To find the optimal parameter $w$ which best approximates the true state value function, we minimize the objective function. 
        In particuluar, the objective function is given as
        \begin{equation}
            J(w) = \E[(v_\pi(S) - \hat{v}(S, w))^2],
        \end{equation}
        where $S$ is a random variable of state space. Then what is the distribution of a random variable $S$?
        We can use \emph{stationary distribution} of Markov Process. Stationary distribution of a Markov Process describes the long-term behavior of a Markov decision process.
        That is, the probability of agent being located at certain state after a long execution can be described by this distribution.

        Let $d_\pi(s)$ be the stationary distribution of the Markov process under policy $\pi$. Then, the objective functon can be written as
        \begin{equation}
            J(w) = \sum_{s \in \mathcal{S}}d_\pi(s)(v_\pi(s) - \hat{v}(s, w))^2 \label{eq:stationaryDist}
        \end{equation}
        TODO: Why stationary distribution? + Tsitsiklis and van Roy + Chapter 6

\section{Deep Q-Learning}
    \subsection{Main Idea}
        The \textbf{Deep Q-Learning} \cite{Mnih2013} is one of the earliest and successful deep reinforcement learning algorithms.
        Suprisingly, DQN algorithm are not guaranteed to converge, since the deadly triad perfectly fits into the algorithms' condition:
        it bootstraps, it is off-policy, and it uses function approximation. But there are practical methods for calming down this unstability.
        But still note that the algorithm is not guaranteed to converge.

    \subsection{Objective Function}
        The objective function which DQN minimizes is given as following:
        \begin{equation}
            J = \E\Big[\Big( R + \gamma \max_{a \in \mathcal{A}(S')}{\hat{q}(S',a,w)} - \hat{q}(S, A, w) \Big)^2\Big],
        \end{equation}
        where $\hat{q}$ is the approximation function of action values
        and $(S, A, R, S')$ are random variables denoting state, action, reward and next state, resepectively.
        Note that the squared term is a Bellman optimality error.

        Now, we can caculate the gradient and minimize the loss with gradient descents... but take a look. Unfortunately, the
        max operation is not differentiable. To resolve this, we freeze the parameter of a network inside the max operation for a short time. Specifically,
        we introduce two networks: \emph{target network and main network}. Using the main network $\hat{q}(s, a, w)$ and $\hat{q}(s, a, w_T)$,
        we can re-write the objective function as following:
        \begin{equation}
            J = \E\Big[\Big( R + \gamma \max_{a \in \mathcal{A}(S')}{\hat{q}(S',a,w_T)} - \hat{q}(S, A, w) \Big)^2\Big], \label{eq:DQNObjective}
        \end{equation}
        and the gradient can be computed about $w$. The parameters of target network and main network are synchronized regularly.

        There is one more technique required for DQN to be trained. The \emph{experience replay} collects the $(s, a, r, s')$ pair, and then
        samples the experiences with uniform random distrubution. Why is this technique required?
        
        \begin{itemize}
            \item \cite{Zhao2025} says: Compare equation \ref{eq:DQNObjective} and \ref{eq:stationaryDist}. Since we don't know the probability distribution of $(S, A, R, S')$,
                we just simply set it as uniform distrubution. However, the state-action pairs are are strongly correlated, and not uniformly distributed in single trajectory.
                To break the correlation, we uniformly draw samples from the replay buffer.
        \end{itemize}
        
        Also, since we can reuse the experiences, the sample efficiency naturally increases.
    \subsection{Pseudocode}
        \Algorithm{DQN}
        {Deep Q-Learning}
        {
            \inputitem{$b$}{behavior policy} \\
            \inputitem{$C$}{sync rate}
        }
        {\inputitem{$\hat{q}(s, a, w)$}{approximated action values}}
        {
            store the samples generated by $b$ in a replay buffer $\mathcal{B} = \{(s, a, r, s')\}$\;
            \ForEach{iteration}
            {
                uniformly draw a mini batch $B$ of samples from $\mathcal{B}$\;
                $L \leftarrow 0$\;
                \ForEach{sample $(s, a, r, s') \in B$}
                {
                    $y_T \leftarrow r + \gamma \max_{a \in \mathcal{A}(s')}{\hat{q}(s', a, w_T)}$\;
                    $L \leftarrow L + \Big( y_T - \hat{q}(s, a, w) \Big)^2$\;
                }
                minimize $L$ by updating parameter $w$\;
                set $w_t = w$ every $C$ iterations
            }
            \Return{$\hat{q}$}
            
        }
    \subsection{Implementation Details}
        Since the original DQN algorithm is highly unstable, there are several techniques to stabilize the algorithm.
        \begin{itemize}
            \item Priortized experience replay (PER)
            \item Double DQN
            \item Dueling DQN
            \item Distributional DQN
        \end{itemize}
    \subsection{Pros and Cons}
        Pros:
        \begin{itemize}
            \item 
        \end{itemize}
        Cons:
        \begin{itemize}
            \item Catastrophic forgetting: This is a common problems in neural networks.
            \item No convergence guaranteed: Note that DQN satisfies the deadly triad condition.
                And also, the algorithm does not even use the true gradient, since the target network is freezed for a moment. This is called a \textbf{semi-gradient} method,
                and no convergence is guaranteed for this method. For more information see \cite{Sutton2020}.
        \end{itemize}